% Section 3.2, from lectures/L03/README.md: the goals, the evaluation questions, and the next
% lecture.

\section{Review}
\label{c3:sec:review}

\needspace{6\baselineskip}
\subsection*{What you should be able to do now}
After this chapter, you should:
\begin{itemize}
  \item Explain, in your own words, why interrupts exist and what problem they solve compared to
    polling.
  \item Be able to configure an external interrupt (INT0 or INT1) for a specific edge and write its
    ISR.
  \item Be able to configure a Pin Change Interrupt on an arbitrary pin, including working out the
    correct \code{PCINTx} number.
  \item Understand why PCI interrupts fire on both edges, and how to filter for the one you care
    about.
\end{itemize}

\needspace{6\baselineskip}
\subsection*{Questions to test yourself}
You should be able to answer:
\begin{enumerate}
  \item What does \code{EICRA} configure that \code{EIMSK} doesn't, and vice versa?
  \item Why do multiple Pin Change Interrupts on the same port share a single ISR, and what does
    that mean for how you write it?
  \item A PCI-based button handler toggles an LED on both press \emph{and} release. What line is
    missing, and why?
\end{enumerate}

\needspace{6\baselineskip}
\subsection*{Looking ahead}
Chapter~\ref{c4:ch} is about analog-to-digital conversion:
\begin{itemize}
  \item Analog-to-digital conversion (ADC): reading a continuous signal instead of a simple
    high/low.
  \item A worked example combining what you already know (PCI, bitwise register I/O) with the ADC.
\end{itemize}
